\documentclass[12pt,a4paper]{article}

% Packages for formatting and graphics
\usepackage[utf8]{inputenc}
\usepackage[T1]{fontenc}
\usepackage[english]{babel}
\usepackage{graphicx}
\usepackage{hyperref}
\usepackage{geometry}
\usepackage{booktabs}
\usepackage{longtable}
\usepackage{listings}
\usepackage{xcolor}

% Page geometry
\geometry{margin=1in}

% Code listing style
\lstset{
    basicstyle=\ttfamily\small,
    breaklines=true,
    frame=single,
    backgroundcolor=\color{gray!5},
    keywordstyle=\color{blue},
    commentstyle=\color{green!50!black},
    stringstyle=\color{red},
    captionpos=b
}

\title{Assignment 2: Requirements, Architecture, and MVP Deployment \\ \large Poetry: A Conversational Recommender System}
\author{Zakhar (solo Developer)}
\date{\today}

\begin{document}

\maketitle

\tableofcontents
\newpage

\section{Use Case Diagram Description}
The Poetry Recommender system is designed to provide a seamless interaction between the user and a rich database of classical poetry. The primary actors in the system are the \textbf{User} and the \textbf{Administrator}.

\subsection{Primary Use Cases}
\begin{itemize}
    \item \textbf{Register/Set Preferences}: The User initiates contact with the Telegram bot and specifies their language preference (English or Russian).
    \item \textbf{Get Recommendation}: The system analyzes the user's profile and context to suggest a poem.
    \item \textbf{Learn/Memorize Poem}: The User adds a poem to their personal learning list, triggering the SM-2 spaced repetition scheduler.
    \item \textbf{Voice Recitation Check}: The User sends a voice recording of themselves reciting a poem. The system uses Vosk STT to evaluate the accuracy and provides feedback.
    \item \textbf{View Progress}: The User checks their learning statistics, XP, level, and current streaks.
    \item \textbf{Seed Database (Admin)}: The Administrator populates the system with high-quality poetry data.
\end{itemize}

\begin{figure}[h]
    \centering
    % TODO: Replace with your Use Case Diagram image
    \fbox{Placeholder for Use Case Diagram Image}
    \caption{Use Case Diagram for Poetry Recommender System}
    \label{fig:usecase}
\end{figure}

\newpage
\section{User Stories}
Below are 15 high-quality user stories following the strict Agile format, prioritized and estimated with Story Points (SP).

\begin{longtable}{@{}p{0.1\textwidth}p{0.65\textwidth}p{0.15\textwidth}@{}}
\toprule
\textbf{ID} & \textbf{User Story (As a... I want to... So that...)} & \textbf{Points (SP)} \\ \midrule
US001 & As a new user, I want to start the bot using /start, so that I can register and set my language preferences. & 2 \\
US002 & As a user, I want to receive a personalized poem recommendation, so that I can discover new poetry I like. & 8 \\
US003 & As a user, I want to add a recommended poem to my learning list, so that I can start memorizing it. & 3 \\
US004 & As a student, I want to use the SM-2 algorithm for learning, so that I review the poems at optimal intervals. & 13 \\
US005 & As a user, I want to recite a poem via voice message, so that the bot can automatically check my accuracy. & 13 \\
US006 & As a user, I want to see an accuracy score after reciting a poem, so that I know which parts I missed. & 5 \\
US007 & As a user, I want to manually rate my recall (0-5), so that I can track my progress without voice recording. & 3 \\
US008 & As a user, I want to receive daily reminders for due reviews, so that I don't break my learning streak. & 5 \\
US009 & As a competitive user, I want to earn XP and level up for every review, so that I feel motivated to keep learning. & 5 \\
US010 & As a user, I want to view my progress dashboard, so that I can see my memorization statistics. & 5 \\
US011 & As a user, I want to maintain a "streak" of consecutive learning days, so that I can stay disciplined. & 3 \\
US012 & As a user, I want to search for poems by author or title, so that I can find specific works to learn. & 5 \\
US013 & As a developer, I want to deploy the system using Docker Compose, so that the environment is reproducible. & 5 \\
US014 & As an admin, I want to seed the database with initial poems, so that the system is ready for users. & 3 \\
US015 & As a user, I want my data stored securely in a PostgreSQL database, so that my progress is persistent. & 5 \\
\bottomrule
\end{longtable}

\newpage
\section{Product Backlog \& Task Management}
The project utilizes a \textbf{DEEP} backlog (Detailed appropriately, Emergent, Estimated, and Prioritized) managed via GitHub Projects.

\subsection{GitHub Kanban Management}
Tasks are tracked using a Kanban board with the following columns:
\begin{itemize}
    \item \textbf{Backlog}: All future features and user stories.
    \item \textbf{To Do}: Tasks selected for the current iteration.
    \item \textbf{In Progress}: Tasks currently being worked on.
    \item \textbf{Done}: Completed and verified tasks.
\end{itemize}

Each task includes clear acceptance criteria and is linked to relevant user stories. Pull Requests are used for all code changes, ensuring that the main branch remains stable.

\textbf{Link to GitHub Kanban}: \href{TODO: Placeholder for GitHub Kanban link}{GitHub Project Board}

\newpage
\section{Data Modeling \& Architecture}

\subsection{Entity-Relationship Diagram (ERD)}
The data model focuses on the relationships between users, poems, and their memorization progress.

\begin{lstlisting}[language=PlantUML, caption=PlantUML ERD Code]
@startuml
entity "User" as user {
  * id : UUID <<PK>>
  --
  * telegram_id : BigInt <<unique>>
  username : String
  first_name : String
  language_pref : String
  preferences : JSON
  xp : Integer
  level : Integer
  streak : Integer
  created_at : Timestamp
  updated_at : Timestamp
}

entity "Poem" as poem {
  * id : UUID <<PK>>
  --
  * title : String
  * author : String
  * text : Text
  * language : String
  difficulty : Float
  lines_count : Integer
  themes : JSON
  era : String
  embedding : Vector(384)
}

entity "Memorization" as memorization {
  * id : UUID <<PK>>
  --
  * user_id : UUID <<FK>>
  * poem_id : UUID <<FK>>
  status : Enum (new, learning, reviewing, memorized)
  ease_factor : Float
  interval_days : Integer
  repetitions : Integer
  next_review_at : Timestamp
  last_reviewed_at : Timestamp
  score_history : JSON
  recommended_at : Timestamp
}

user ||--o{ memorization
poem ||--o{ memorization
@enduml
\end{lstlisting}

\subsection{Tech Stack Description}
\begin{itemize}
    \item \textbf{FastAPI}: High-performance Python web framework used for the backend API.
    \item \textbf{PostgreSQL 16 + pgvector}: Relational database for persistent storage, with vector extensions for semantic poem recommendations.
    \item \textbf{aiogram 3.x}: Modern, asynchronous Telegram Bot framework.
    \item \textbf{Vosk}: Efficient, offline Speech-to-Text engine for evaluating voice recordings.
    \item \textbf{Docker}: Containerization for consistent deployment across environments.
\end{itemize}

\subsection{API Mocking and Documentation}
The API is documented using FastAPI's built-in Swagger UI. For external integration and testing, Postman mocks are used.

\textbf{Postman Link}: \href{TODO: Placeholder for Postman link}{Postman Collection} \\
\textbf{Figma Prototype}: \href{TODO: Placeholder for Figma link}{UI Prototype}

\newpage
\section{MVP v0 Deployment \& Customer Feedback}

\subsection{Deployment Summary}
The initial MVP (v0) was deployed on a Linux VPS using \textbf{Docker Compose}. The deployment process ensures that the database, backend, and bot services are isolated and easily manageable.

\begin{lstlisting}[language=bash]
# Deployment command
docker-compose up -d --build
\end{lstlisting}

\subsection{Customer Review and Phase Definition}
A review session with the "customer" (early user) revealed high satisfaction with the voice check feature. Based on this feedback, the project phases were refined:
\begin{itemize}
    \item \textbf{Phase 1}: Core recommendation and memorization engine (Completed).
    \item \textbf{Phase 2}: Gamification and Leaderboard implementation (Current).
    \item \textbf{Phase 3}: Advanced UI/UX and Mobile responsiveness (Future).
\end{itemize}

\subsection{MVP Roadmap Update}
Based on feedback, the roadmap was updated to include:
\begin{itemize}
    \item Comprehensive leaderboard for competitive motivation.
    \item Stanza-by-stanza learning mode for longer poems.
    \item Enhanced progress visualization in the bot interface.
\end{itemize}

\section{Activity Tracking}
All development activities, including time spent on requirements, design, and coding, are tracked in the Activity Tracking Sheet.

\textbf{Link}: \href{TODO: Placeholder for Activity Tracking Sheet link}{Activity Sheet}

\end{document}
